\chapter{Mixed Set D: Full Investigations}

The first three sets drill one skill per problem. This set does the
opposite. Each problem hands you a single function of two variables and
runs it through the whole chain of questions that a function of two
variables invites: where it lives, what values it takes, where it can be
differentiated, what its level curves look like, where its gradient
points, what its Hessian says, whether it has critical points, which of
its contour sets are convex, and finally a checklist of properties, each
to be decided \emph{with a reason}.

None of the steps is new. What is new is that they feed each other: the
domain decides where the derivative may be taken, the Hessian decides
concavity, concavity decides quasiconcavity, the domain and the range
decide whether a maximum exists. A slip early in the chain propagates,
and a good investigation catches it, because the later answers stop
fitting together. Target time: twenty to twenty-five minutes per
function. The solutions tag every step with the volume and chapter whose
rule it uses, so that a weak link can be traced back and repaired.

\section{The Investigation Chain}

\begin{rulebox}{The chain, in order}
\begin{enumerate}[label=\arabic*.,leftmargin=1.6em,itemsep=0.3ex]
  \item \textbf{Domain.} Collect the obstructions (roots need a
        nonnegative radicand, logarithms a positive argument,
        denominators must not vanish), sketch the set, and say whether it
        is closed, bounded and convex. (Vol.~2, Ch.~1 and Ch.~4;
        Vol.~5, Ch.~10; Vol.~7, Ch.~2.)
  \item \textbf{Range.} Find the smallest and largest values, or show that
        there are none, and use continuity on a connected domain to fill
        in everything between. (Vol.~2, Ch.~1; Vol.~5, Ch.~9.)
  \item \textbf{Derivative.} Compute $f_{x}$, $f_{y}$ and say where they
        exist and are continuous: on that open set $f$ is differentiable
        (continuous partials are sufficient, not necessary). Where a
        partial fails to exist, $f$ is not differentiable. Roots are the
        usual offenders on the boundary. (Vol.~5, Ch.~10.)
  \item \textbf{Level curves.} Solve $f(x,y)=c$ for one variable and
        sketch several values of $c$. (Vol.~5, Ch.~9.)
  \item \textbf{Gradient at a point.} Evaluate $\grad f$, draw it, check
        that it is orthogonal to the level curve through the point, and
        read it: it points in the direction of steepest ascent, and its
        norm $\lVert\grad f\rVert$ is the slope in that direction.
        (Vol.~5, Ch.~10 and Ch.~13.)
  \item \textbf{Hessian.} Second partials, Young's theorem, and the
        definiteness by Sylvester's criterion. (Vol.~5, Ch.~11.)
  \item \textbf{Critical points.} Solve $\grad f=\mathbf{0}$ \emph{inside
        the domain}, or explain why there is no solution; classify with the
        Hessian. (Vol.~7, Ch.~5.)
  \item \textbf{Contour sets.} Mark an upper set $\{f\ge c\}$ and a lower
        set $\{f\le c\}$ and decide whether each is convex.
        (Vol.~5, Ch.~9; Vol.~7, Ch.~2.)
  \item \textbf{Checklist.} Concave, convex, quasiconcave, quasiconvex,
        has a maximum, has a minimum, increasing, decreasing, epigraph
        convex: each answered with a proof or a counterexample.
        (Vol.~7, Ch.~1 and Ch.~3.)
\end{enumerate}
\end{rulebox}

\begin{rulebox}{Three definitions the checklist needs}
\begin{itemize}[leftmargin=1.6em,itemsep=0.3ex]
  \item \textbf{Monotone in several variables.} Order points componentwise:
        $(x,y)\le(x',y')$ means $x\le x'$ and $y\le y'$. Then $f$ is
        \emph{increasing} \de{monoton wachsend} if
        $(x,y)\le(x',y')\Rightarrow f(x,y)\le f(x',y')$, and
        \emph{decreasing} if the inequality between the values is reversed.
        On a convex domain with interior points, a continuous $f$ with
        $f_{x}\ge0$ and $f_{y}\ge0$ at every interior point is increasing.
        To refute it, exhibit two ordered points whose values are in the
        wrong order.
  \item \textbf{Epigraph and hypograph.} The \emph{epigraph}
        \de{Epigraph} of $f$ is the set of points on or above its graph,
        $\{(x,y,z): (x,y)\in D,\ z\ge f(x,y)\}$; the \emph{hypograph} is
        the set on or below it. On a convex domain, the epigraph is a convex
        set exactly when $f$ is convex, and the hypograph exactly when $f$
        is concave. So the epigraph question is the convexity question
        asked again, and the two answers must agree.
  \item \textbf{Existence of extrema.} A continuous function on a closed
        and bounded domain attains a maximum and a minimum (Weierstrass,
        Vol.~7, Ch.~1). On an unbounded or open domain, look at the range
        first: if $f$ grows without bound there is no maximum, and if the
        infimum is approached but never reached there is no minimum.
        Extrema on the boundary of the domain are not found by
        $\grad f=\mathbf{0}$, and are often points where $f$ is not even
        differentiable.
\end{itemize}
\end{rulebox}

\begin{rulebox}{Shortcuts that keep the checklist consistent}
\begin{itemize}[leftmargin=1.6em,itemsep=0.3ex]
  \item Concave $\Rightarrow$ quasiconcave; convex $\Rightarrow$
        quasiconvex. A function that is both concave and convex is affine.
  \item A strictly increasing transform $F$ ($\ln$ of a positive $f$,
        $\sqrt{\cdot}$ or a positive power of a nonnegative $f$) only
        relabels the contour sets, $\{F(f)\ge F(c)\}=\{f\ge c\}$, so it
        preserves quasiconcavity and quasiconvexity but not necessarily
        concavity.
  \item If $g$ is concave and $F$ is increasing and concave, $F\circ g$ is
        concave: this is how $\sqrt{\text{concave}}$ and
        $\ln(\text{concave})$ are handled without a Hessian.
  \item One straight line in the domain along which $f$ is strictly convex
        refutes concavity; one along which it is strictly concave refutes
        convexity. Two points in a contour set whose midpoint is not in it
        refute quasiconcavity or quasiconvexity.
  \item Cobb--Douglas $x^{a}y^{b}$ ($a,b>0$) on $\R^{2}_{+}$: always
        quasiconcave; concave iff $a+b\le1$; never convex.
        (Vol.~7, Ch.~3.)
\end{itemize}
\end{rulebox}

\section{Worked Example}

\begin{worked}{Investigate $f(x,y)=\sqrt{8-2x^{2}-y^{2}}$ completely: domain, range, derivative, level curves, gradient at $(1,2)$, Hessian, critical points, contour sets, and the full property checklist}
\wstep{Step 1. Domain (Vol.~2, Ch.~1; Vol.~2, Ch.~4).}
The only obstruction is the root: $8-2x^{2}-y^{2}\ge0$, that is
\[
  D=\{(x,y): 2x^{2}+y^{2}\le8\}
   =\Bigl\{\tfrac{x^{2}}{4}+\tfrac{y^{2}}{8}\le1\Bigr\},
\]
the filled ellipse with semi-axes $2$ along $x$ and $2\sqrt2$ along $y$,
boundary included. It is \textbf{closed} (the inequality is weak),
\textbf{bounded}, and \textbf{convex}: it is a lower contour set of the
convex function $2x^{2}+y^{2}$ (Vol.~7, Ch.~2).
\wtrap{Writing $2x^{2}+y^{2}<8$ loses the boundary. The root of zero is
zero, so the boundary ellipse belongs to the domain, and it will carry the
minimum.}

\wstep{Step 2. Range (Vol.~2, Ch.~1).}
The radicand is largest at the origin, where it equals $8$, and smallest
($0$) on the boundary. Since $\sqrt{\cdot}$ is increasing,
$0\le f\le\sqrt8=2\sqrt2$. $f$ is continuous on the connected set $D$, so
every value in between is taken:
\[
  f(D)=[\,0,\,2\sqrt2\,].
\]

\wstep{Step 3. Derivative, and where it exists (Vol.~5, Ch.~10).}
By the chain rule,
\[
  f_{x}=\frac{-4x}{2\sqrt{8-2x^{2}-y^{2}}}=-\frac{2x}{f},\qquad
  f_{y}=\frac{-2y}{2\sqrt{8-2x^{2}-y^{2}}}=-\frac{y}{f}.
\]
Both are continuous wherever $f>0$, so $f$ is differentiable on the
\textbf{interior} $2x^{2}+y^{2}<8$. On the boundary the formulas divide by
zero, and the difference quotient confirms the failure: at $(2,0)$,
$\bigl(f(2-h,0)-f(2,0)\bigr)/h=\sqrt{8h-2h^{2}}/h\to\infty$ as
$h\to0^{+}$. The surface meets the plane $z=0$ vertically, like the edge
of a dome.

\wstep{Step 4. Level curves (Vol.~5, Ch.~9).}
For $0\le c\le2\sqrt2$, $f=c$ means $2x^{2}+y^{2}=8-c^{2}$: concentric
ellipses with the same shape as the boundary, shrinking to the origin as
$c\to2\sqrt2$. The values $c=0,1,\sqrt2,2,\sqrt7$ give
$2x^{2}+y^{2}=8,7,6,4,1$ (left panel below). The level $c=0$ is the
boundary of the domain; the level $c=2\sqrt2$ is the single point
$(0,0)$.

\wstep{Step 5. Gradient at $(1,2)$ (Vol.~5, Ch.~10).}
Here $f(1,2)=\sqrt{8-2-4}=\sqrt2$, so
\[
  \grad f(1,2)=\begin{pmatrix}-2/\sqrt2\\-2/\sqrt2\end{pmatrix}
  =\begin{pmatrix}-\sqrt2\\-\sqrt2\end{pmatrix},
  \qquad
  \lVert\grad f(1,2)\rVert=\sqrt{2+2}=2 .
\]
\emph{Orthogonality.} The level curve through $(1,2)$ is
$2x^{2}+y^{2}=6$; implicit differentiation gives
$y'=-4x/(2y)=-1$ there, so a tangent direction is $(1,-1)$, and
$(-\sqrt2)(1)+(-\sqrt2)(-1)=0$ \checkmark.
\emph{Direction.} The gradient points into the ellipse, towards higher
levels: the direction of steepest ascent. \emph{Norm.} $2$ is the slope of
the surface in that direction, the largest directional derivative at
$(1,2)$ (Vol.~5, Ch.~13). In every direction along the level curve the
slope is $0$.

\begin{gallery}
\begin{tikzpicture}
  \begin{axis}[contouraxis, xmin=-3.2, xmax=3.2, ymin=-3.2, ymax=3.2,
      xtick={-2,-1,1,2}, ytick={-2,-1,1,2},
      title={domain, levels, $\grad f(1,2)$}]
    \addplot[name path=top, draw=none, domain=-2:2]
      {sqrt(max(8-2*x^2,0))};
    \addplot[name path=bot, draw=none, domain=-2:2]
      {-sqrt(max(8-2*x^2,0))};
    \region{top}{bot}{}
    \addplot[accent, line width=1pt, domain=0:360, samples=121]
      ({2*cos(x)},{sqrt(8)*sin(x)});
    \addplot[domain=0:360, samples=121]({sqrt(3.5)*cos(x)},{sqrt(7)*sin(x)});
    \addplot[domain=0:360, samples=121]({sqrt(3)*cos(x)},{sqrt(6)*sin(x)});
    \addplot[domain=0:360, samples=121]({sqrt(2)*cos(x)},{2*sin(x)});
    \addplot[domain=0:360, samples=121]({sqrt(0.5)*cos(x)},{sin(x)});
    \draw[dashed, gray!70!black, line width=0.5pt]
      (axis cs:0.1,2.9) -- (axis cs:1.9,1.1);
    \draw[-{Stealth[length=6pt,width=4pt]}, accent, line width=1pt]
      (axis cs:1,2) -- (axis cs:-0.414,0.586);
    \addplot[only marks, mark=*, mark size=1.2pt, black, forget plot]
      coordinates {(1,2) (0,0)};
  \end{axis}
\end{tikzpicture}\hfill
\begin{tikzpicture}
  \begin{axis}[contouraxis, xmin=-3.2, xmax=3.2, ymin=-3.2, ymax=3.2,
      xtick={-2,-1,1,2}, ytick={-2,-1,1,2},
      title={upper set $\{f\ge2\}$ (shaded)}]
    \addplot[name path=itop, draw=none, domain=-1.4142:1.4142]
      {sqrt(max(4-2*x^2,0))};
    \addplot[name path=ibot, draw=none, domain=-1.4142:1.4142]
      {-sqrt(max(4-2*x^2,0))};
    \region{itop}{ibot}{}
    \addplot[accent, line width=1pt, domain=0:360, samples=121]
      ({2*cos(x)},{sqrt(8)*sin(x)});
    \addplot[domain=0:360, samples=121]({sqrt(2)*cos(x)},{2*sin(x)});
    \draw[gray!70!black, line width=0.5pt] (axis cs:0,2) -- (axis cs:0,-2);
    \addplot[only marks, mark=*, mark size=1.2pt, black, forget plot]
      coordinates {(0,2) (0,-2)};
    \addplot[only marks, mark=o, mark size=1.6pt, black, forget plot]
      coordinates {(0,0)};
  \end{axis}
\end{tikzpicture}
\gcap{Left: the domain (shaded), its boundary $f=0$ (bold), the levels
$f=1,\sqrt2,2,\sqrt7$, and the gradient at $(1,2)$, orthogonal to the
dashed tangent of its level curve. Right: the upper set $\{f\ge2\}$ is a
filled ellipse; the lower set $\{f\le2\}$ is the ring between the two
curves, which contains $(0,\pm2)$ but not their midpoint.}
\end{gallery}

\wstep{Step 6. Hessian (Vol.~5, Ch.~11).}
Differentiate $f_{x}=-2xf^{-1}$ and $f_{y}=-yf^{-1}$ once more, using
$f_{x}$, $f_{y}$ from Step 3:
\[
  f_{xx}=-\frac{2}{f}+\frac{2x f_{x}}{f^{2}}
        =-\frac{2f^{2}+4x^{2}}{f^{3}}
        =-\frac{2(8-y^{2})}{f^{3}},\qquad
  f_{yy}=-\frac{1}{f}+\frac{y f_{y}}{f^{2}}
        =-\frac{8-2x^{2}}{f^{3}},
\]
\[
  f_{xy}=\frac{2x f_{y}}{f^{2}}=-\frac{2xy}{f^{3}}
  \quad(\text{equal to } f_{yx}\text{ by Young}).
\]
The determinant simplifies remarkably:
\[
  \det\Hess f=\frac{2(8-y^{2})(8-2x^{2})-4x^{2}y^{2}}{f^{6}}
  =\frac{16(8-2x^{2}-y^{2})}{f^{6}}=\frac{16}{f^{4}}>0 .
\]
In the interior $y^{2}<8$, so $f_{xx}<0$: the Hessian is
\textbf{negative definite} at every interior point. At $(1,2)$, with
$f^{3}=2\sqrt2$:
\[
  \Hess f(1,2)=\begin{pmatrix}-2\sqrt2&-\sqrt2\\[2pt]-\sqrt2&-\tfrac32\sqrt2
  \end{pmatrix},\qquad \det=6-2=4=\frac{16}{(\sqrt2)^{4}}\ \checkmark
\]
\wtrap{Do not expand $\det\Hess f$ with the roots still in it. Writing
every entry over the common denominator $f^{3}$ and replacing $f^{2}$ by
$8-2x^{2}-y^{2}$ is what makes the sign visible.}

\wstep{Step 7. Critical points (Vol.~7, Ch.~5).}
$f_{x}=0\iff x=0$ and $f_{y}=0\iff y=0$, so the only candidate is
$(0,0)$, which lies in the interior. There
$\Hess f(0,0)=\begin{pmatrix}-1/\sqrt2&0\\0&-1/(2\sqrt2)\end{pmatrix}$ is
negative definite: a strict local maximum, with value $2\sqrt2$. The
boundary points are not critical points; $f$ is not differentiable there.

\wstep{Step 8. Contour sets (Vol.~5, Ch.~9; Vol.~7, Ch.~2).}
Take $c=2$. The upper set $\{f\ge2\}=\{2x^{2}+y^{2}\le4\}$ is a filled
ellipse: \textbf{convex}. The lower set
$\{f\le2\}=\{4\le2x^{2}+y^{2}\le8\}$ is an elliptical ring: it contains
$(0,2)$ and $(0,-2)$ but not their midpoint $(0,0)$, so it is
\textbf{not convex} (right panel).

\wstep{Step 9. The checklist, each with its reason.}
\begin{itemize}[leftmargin=1.6em,itemsep=0.3ex]
  \item \emph{Concave: yes.} The radicand $8-2x^{2}-y^{2}$ is concave and
        $\sqrt{\cdot}$ is increasing and concave, so the composition is
        concave on $D$ (Vol.~7, Ch.~3). Step 6 confirms it, even strictly,
        on the interior.
  \item \emph{Convex: no.} The Hessian is negative definite; concretely,
        $f(0,0)=2\sqrt2$ lies above the chord value
        $\tfrac12\bigl(f(-2,0)+f(2,0)\bigr)=0$.
  \item \emph{Quasiconcave: yes}, because concave; every upper set is a
        filled ellipse (Step 8).
  \item \emph{Quasiconvex: no.} The lower set $\{f\le2\}$ is not convex
        (Step 8).
  \item \emph{Maximum: yes}, $2\sqrt2$ at $(0,0)$, unique; global because
        $f$ is concave and $(0,0)$ is a critical point (Vol.~7, Ch.~5).
  \item \emph{Minimum: yes}, $0$, attained at \emph{every} point of the
        boundary ellipse. Weierstrass guarantees it (continuous on a closed
        bounded set); the first-order condition could never find it, since
        $f$ is not differentiable there.
  \item \emph{Increasing: no.} $(0,0)\le(1,0)$ but
        $f(1,0)=\sqrt6<2\sqrt2=f(0,0)$.
  \item \emph{Decreasing: no.} $(-1,0)\le(0,0)$ but
        $f(-1,0)=\sqrt6<f(0,0)$.
  \item \emph{Epigraph convex: no}, since $f$ is not convex. Its
        hypograph, the solid dome, is convex, matching concavity.
\end{itemize}

\wstep{Step 10. Do the answers fit together?}
Concave with an interior critical point gives a global maximum
\checkmark. Not quasiconvex, yet concave: a dome has convex upper sets and
ring-shaped lower sets \checkmark. The range $[0,2\sqrt2]$ matches the
maximum and minimum found \checkmark. Neither monotone, as one expects of
a function that peaks in the middle of its domain \checkmark.
\end{worked}

\section{Problems}

Each problem is one investigation. Work every part, and give a reason for
every item of the checklist.

\begin{problems}
  \item Investigate $f(x,y)=\sqrt{6-x-y^{2}}$.
        \begin{enumerate}[label=(\alph*),leftmargin=1.8em,itemsep=0.2ex]
          \item Give and sketch the domain. Is it closed, bounded, convex?
                What values does $f$ take?
          \item Compute $f_{x}$ and $f_{y}$. Is $f$ differentiable on its
                whole domain?
          \item Describe and sketch the level curves.
          \item Compute $\grad f(1,1)$, draw it, show that it is orthogonal
                to the level curve through $(1,1)$, and state what its
                direction and its norm mean.
          \item Compute the Hessian and its definiteness.
          \item Does $f$ have critical points?
          \item Mark an upper and a lower contour set at level $2$. Which is
                convex?
          \item Decide, with reasons: concave, convex, quasiconcave,
                quasiconvex, has a maximum, has a minimum, increasing,
                decreasing, epigraph convex.
        \end{enumerate}
  \item Investigate the utility function
        $u:\R^{2}_{++}\to\R$, $u(x,y)=\ln\bigl(x^{1/3}y^{1/2}\bigr)$, where
        $\R^{2}_{++}=\{(x,y):x>0,\ y>0\}$.
        \begin{enumerate}[label=(\alph*),leftmargin=1.8em,itemsep=0.2ex]
          \item Rewrite $u$ as a sum. Describe the domain (closed? bounded?
                convex?) and the range.
          \item Describe and sketch the level curves. How do they compare
                with those of $x^{1/3}y^{1/2}$?
          \item Compute $\grad u(1,1)$, draw it, check orthogonality to the
                level curve, and interpret direction and norm.
          \item Compute the Hessian. Show that both marginal utilities are
                diminishing, and give the definiteness.
          \item Does $u$ have critical points?
          \item Mark an upper and a lower contour set at level $0$. Which is
                convex?
          \item Decide, with reasons: concave, convex, quasiconcave,
                quasiconvex, has a maximum, has a minimum, increasing,
                decreasing, epigraph convex.
        \end{enumerate}
  \item Investigate the production function $f:\R^{2}_{+}\to\R$,
        $f(x,y)=x^{3/4}y^{1/2}$.
        \begin{enumerate}[label=(\alph*),leftmargin=1.8em,itemsep=0.2ex]
          \item Name this type of function and its returns to scale. Give
                the domain (closed? convex?) and the range.
          \item Compute the derivative and say where it exists. Show that
                both marginal products are diminishing.
          \item Compute the Hessian and its definiteness. Does diminishing
                marginal productivity make $f$ concave?
          \item Sketch the level curves, and draw $\grad f(1,1)$; check
                orthogonality and interpret direction and norm.
          \item Does $f$ have critical points? Where does its minimum come
                from?
          \item Mark an upper and a lower contour set at level $1$. Which is
                convex?
          \item Decide, with reasons: concave, convex, quasiconcave,
                quasiconvex, has a maximum, has a minimum, increasing,
                decreasing, epigraph convex.
        \end{enumerate}
  \item Investigate $f:\R^{2}_{+}\to\R$, $f(x,y)=x^{1/3}y^{2/3}$.
        \begin{enumerate}[label=(\alph*),leftmargin=1.8em,itemsep=0.2ex]
          \item Name the function and its returns to scale; give domain and
                range.
          \item Compute the derivative and the Hessian, and the
                definiteness. Which direction does the Hessian annihilate?
          \item Sketch the level curves and draw $\grad f(1,8)$; check
                orthogonality, interpret direction and norm, and verify
                Euler's theorem at $(1,8)$.
          \item Does $f$ have critical points?
          \item Mark an upper and a lower contour set at level $4$. Which is
                convex?
          \item Decide, with reasons: concave, strictly concave, convex,
                quasiconcave, quasiconvex, has a maximum, has a minimum,
                increasing, decreasing, epigraph convex.
        \end{enumerate}
  \item Investigate $f(x,y)=\ln\bigl(4+x^{2}+y^{2}\bigr)$.
        \begin{enumerate}[label=(\alph*),leftmargin=1.8em,itemsep=0.2ex]
          \item Give the domain and the range.
          \item Compute the gradient; where is $f$ differentiable?
          \item Describe the level curves, compute $\grad f(1,1)$, and check
                that it is orthogonal to the level curve through $(1,1)$.
          \item Compute the Hessian, its determinant in general, and its
                definiteness at $(0,0)$, $(1,1)$ and $(2,2)$.
          \item Find and classify the critical points.
          \item Mark the upper and lower contour sets at level $\ln8$.
                Which is convex?
          \item Decide, with reasons: concave, convex, quasiconcave,
                quasiconvex, has a maximum, has a minimum, increasing,
                decreasing, epigraph convex. On which disk centred at the
                origin is $f$ convex?
        \end{enumerate}
\end{problems}

\section{Solutions}

\begin{solutions}
  \item \textbf{Root of a concave expression on an unbounded closed domain.}
        \begin{enumerate}[label=(\alph*),leftmargin=1.8em,itemsep=0.4ex]
          \item \emph{Domain (Vol.~2, Ch.~1 and Ch.~4).} Radicand
                $\ge0$: $D=\{x\le6-y^{2}\}$, the region on and to the left of
                the parabola $x=6-y^{2}$ (vertex $(6,0)$, opening to the
                left). Closed (weak inequality), unbounded (it contains the
                whole ray $y=0$, $x\le6$), and convex: it is
                the upper set $\{g\ge0\}$ of the concave function
                $g=6-x-y^{2}$ (Vol.~7, Ch.~2). \emph{Range
                (Vol.~2, Ch.~1):} $f=0$ on the parabola and
                $f(x,0)=\sqrt{6-x}\to\infty$ as $x\to-\infty$; by
                continuity $f(D)=[0,\infty)$.
          \item \emph{Derivative (Vol.~5, Ch.~10).}
                $f_{x}=\dfrac{-1}{2\sqrt{6-x-y^{2}}}=-\dfrac{1}{2f}$,
                \;$f_{y}=\dfrac{-2y}{2f}=-\dfrac{y}{f}$, continuous where
                $f>0$. So $f$ is differentiable on the interior
                $x<6-y^{2}$ only. On the parabola the partials blow up:
                at $(6,0)$, $\bigl(f(6-h,0)-0\bigr)/h=\sqrt h/h\to\infty$.
          \item \emph{Level curves (Vol.~5, Ch.~9).} $f=c\ge0$ gives
                $x=6-c^{2}-y^{2}$: copies of the boundary parabola, shifted
                left by $c^{2}$. For $c=0,1,2,3$ the vertices sit at
                $x=6,5,2,-3$; higher levels lie further left.
          \item \emph{Gradient (Vol.~5, Ch.~10).} $f(1,1)=\sqrt4=2$, so
                \[
                  \grad f(1,1)=\begin{pmatrix}-1/4\\-1/2\end{pmatrix},\qquad
                  \lVert\grad f(1,1)\rVert=\sqrt{\tfrac1{16}+\tfrac14}
                  =\frac{\sqrt5}{4}\approx0.56 .
                \]
                The level curve through $(1,1)$ is $x=2-y^{2}$; there
                $\mathrm{d}x/\mathrm{d}y=-2y=-2$, so
                $\mathrm{d}y/\mathrm{d}x=-\tfrac12$ and $(2,-1)$ is a tangent
                direction. $(-\tfrac14)(2)+(-\tfrac12)(-1)=0$ \checkmark.
                The gradient points left and down, towards the inside of the
                parabolas where the levels rise: steepest ascent, with slope
                $\sqrt5/4$ in that direction (Vol.~5, Ch.~13).
          \item \emph{Hessian (Vol.~5, Ch.~11).} From $f_{x}=-\tfrac12f^{-1}$
                and $f_{y}=-yf^{-1}$:
                \[
                  f_{xx}=\tfrac12f^{-2}f_{x}=-\frac{1}{4f^{3}},\qquad
                  f_{xy}=\tfrac12f^{-2}f_{y}=-\frac{y}{2f^{3}},
                \]
                \[
                  f_{yy}=-\frac1f+\frac{yf_{y}}{f^{2}}
                        =-\frac{f^{2}+y^{2}}{f^{3}}=-\frac{6-x}{f^{3}} .
                \]
                $\det\Hess f=\dfrac{(6-x)-y^{2}}{4f^{6}}
                =\dfrac{f^{2}}{4f^{6}}=\dfrac{1}{4f^{4}}>0$ and
                $f_{xx}<0$: \textbf{negative definite} on the interior. At
                $(1,1)$: $\Hess f=\begin{pmatrix}-1/32&-1/16\\-1/16&-5/8
                \end{pmatrix}$, $\det=\tfrac{5}{256}-\tfrac{1}{256}
                =\tfrac1{64}=\tfrac{1}{4\cdot2^{4}}$ \checkmark.
          \item \emph{Critical points (Vol.~7, Ch.~5).} None:
                $f_{x}=-1/(2f)$ is never $0$. Moving left always raises $f$.
          \item \emph{Contour sets (Vol.~5, Ch.~9; Vol.~7, Ch.~2).}
                $\{f\ge2\}=\{x\le2-y^{2}\}$, the inside of a parabola:
                \textbf{convex}. $\{f\le2\}=\{2-y^{2}\le x\le6-y^{2}\}$, a
                curved band between two parabolas: it contains $(-2,2)$ and
                $(-2,-2)$ (check: $-2\le-2\le2$) but not their midpoint
                $(-2,0)$ (since $2\le-2$ fails). \textbf{Not convex.}
          \item \emph{Checklist (Vol.~7, Ch.~1 and Ch.~3).}
                \emph{Concave: yes}, $\sqrt{\cdot}$ increasing and concave
                applied to the concave $6-x-y^{2}$ (and (e) confirms it).
                \emph{Convex: no}, negative definite Hessian.
                \emph{Quasiconcave: yes}, because concave.
                \emph{Quasiconvex: no}, by the band in (g).
                \emph{Maximum: no}, $f$ is unbounded above.
                \emph{Minimum: yes}, $0$ on the whole parabola $x=6-y^{2}$,
                boundary points where $f$ is not differentiable.
                \emph{Increasing: no}, $f$ falls in $x$: $(0,0)\le(1,0)$
                with $f(1,0)=\sqrt5<\sqrt6=f(0,0)$.
                \emph{Decreasing: no}, $(0,-1)\le(0,0)$ with
                $f(0,-1)=\sqrt5<\sqrt6=f(0,0)$.
                \emph{Epigraph convex: no}, since $f$ is not convex (the
                hypograph is).
        \end{enumerate}
  \item \textbf{Log-Cobb--Douglas.}
        \begin{enumerate}[label=(\alph*),leftmargin=1.8em,itemsep=0.4ex]
          \item \emph{Logarithm rules (Vol.~1, Ch.~2).}
                $u=\tfrac13\ln x+\tfrac12\ln y$. \emph{Domain
                (Vol.~5, Ch.~10):} the open quadrant $\R^{2}_{++}$, open
                (not closed), unbounded, convex. \emph{Range:}
                $u(t,1)=\tfrac13\ln t$ runs through all of $\R$ as $t$
                runs through $(0,\infty)$, so $u(\R^{2}_{++})=\R$.
          \item \emph{Level curves (Vol.~5, Ch.~9).}
                $\tfrac13\ln x+\tfrac12\ln y=c\iff\ln y=2c-\tfrac23\ln x
                \iff y=\mathrm{e}^{2c}x^{-2/3}$: decreasing, convex curves
                asymptotic to both axes, moving outwards as $c$ grows. They
                are exactly the level curves of $x^{1/3}y^{1/2}$, relabelled
                ($u=\ln v$ with $v$ the Cobb--Douglas): an increasing
                transform changes the labels, not the curves.
          \item \emph{Gradient (Vol.~5, Ch.~10).} $u_{x}=\dfrac{1}{3x}$,
                $u_{y}=\dfrac{1}{2y}$, so $\grad u(1,1)=(\tfrac13,\tfrac12)$
                with norm $\sqrt{\tfrac19+\tfrac14}=\tfrac{\sqrt{13}}{6}
                \approx0.60$. The level curve through $(1,1)$ is $u=0$, i.e.
                $y=x^{-2/3}$, with slope $-\tfrac23$ at $x=1$ and tangent
                direction $(3,-2)$; $\tfrac13\cdot3-\tfrac12\cdot2=0$
                \checkmark. The gradient points up and to the right, the
                direction in which utility rises fastest, at rate
                $\sqrt{13}/6$.
          \item \emph{Hessian (Vol.~5, Ch.~11).}
                $\Hess u=\begin{pmatrix}-\frac{1}{3x^{2}}&0\\0&-\frac{1}{2y^{2}}
                \end{pmatrix}$. $u_{xx}<0$ and $u_{yy}<0$: each marginal
                utility $u_{x}=1/(3x)$, $u_{y}=1/(2y)$ falls as its own good
                increases. Diagonal with negative entries: \textbf{negative
                definite} everywhere.
          \item \emph{Critical points (Vol.~7, Ch.~5).} None: $u_{x}=1/(3x)>0$.
          \item \emph{Contour sets.} $u\le0\iff x^{1/3}y^{1/2}\le1\iff
                x^{2}y^{3}\le1$. The upper set $\{u\ge0\}=\{y\ge x^{-2/3}\}$
                lies above a convex curve: \textbf{convex}. The lower set is
                not: $(8,\tfrac14)$ and $(\tfrac18,4)$ both satisfy
                $x^{2}y^{3}=1$, but at the midpoint
                $(\tfrac{65}{16},\tfrac{17}{8})$ the product is far above
                $1$ (already $x^{2}>16$ and $y^{3}>8$).
          \item \emph{Checklist (Vol.~7, Ch.~3).}
                \emph{Concave: yes}, strictly, by the negative definite
                Hessian (or as a positive sum of the concave $\ln x$,
                $\ln y$).
                \emph{Convex: no}, the Hessian is negative definite.
                \emph{Quasiconcave: yes}, because concave.
                \emph{Quasiconvex: no}, by (f).
                \emph{Maximum: no}, $u(t,1)\to\infty$ as $t\to\infty$.
                \emph{Minimum: no}, $u(t,1)\to-\infty$ as $t\to0^{+}$; the
                open domain has no boundary to stop at.
                \emph{Increasing: yes}, strictly: $u_{x}>0$, $u_{y}>0$ on a
                convex domain.
                \emph{Decreasing: no}, as it is strictly increasing.
                \emph{Epigraph convex: no}, $u$ is not convex.
        \end{enumerate}
  \item \textbf{Cobb--Douglas with exponent sum above $1$.}
        \begin{enumerate}[label=(\alph*),leftmargin=1.8em,itemsep=0.4ex]
          \item \emph{Name (Vol.~5, Ch.~14).} A Cobb--Douglas function,
                homogeneous of degree $\tfrac34+\tfrac12=\tfrac54>1$:
                $f(tx,ty)=t^{5/4}f(x,y)$, \textbf{increasing returns to
                scale}. Domain $\R^{2}_{+}$: closed, unbounded, convex.
                Range $[0,\infty)$: $f=0$ on the axes and $f(t,t)=t^{5/4}$
                is unbounded.
          \item \emph{Derivative (Vol.~5, Ch.~10).}
                \[
                  f_{x}=\tfrac34x^{-1/4}y^{1/2},\qquad
                  f_{y}=\tfrac12x^{3/4}y^{-1/2},
                \]
                continuous on the open quadrant $x,y>0$, so $f$ is
                differentiable there. On the axis $x=0$ (with $y>0$) the
                factor $x^{-1/4}$ is infinite, and likewise $y^{-1/2}$ on the
                axis $y=0$ (with $x>0$): the partials do not exist there.
                The origin is the exception: $f$ vanishes on both axes, so
                $f_{x}(0,0)=f_{y}(0,0)=0$, and
                $0\le f(h,k)\le\lVert(h,k)\rVert^{5/4}$ shows that $f$ is
                differentiable there (from inside the quadrant) with zero
                derivative. Second partials:
                $f_{xx}=-\tfrac{3}{16}x^{-5/4}y^{1/2}<0$,
                $f_{yy}=-\tfrac14x^{3/4}y^{-3/2}<0$: both marginal products
                \textbf{diminish}.
          \item \emph{Hessian (Vol.~5, Ch.~11).}
                $f_{xy}=\tfrac38x^{-1/4}y^{-1/2}$, and
                \[
                  \det\Hess f=\tfrac{3}{64}x^{-1/2}y^{-1}
                  -\tfrac{9}{64}x^{-1/2}y^{-1}
                  =-\tfrac{3}{32}\,x^{-1/2}y^{-1}<0 :
                \]
                \textbf{indefinite} at every interior point, so $f$ is
                neither concave nor convex. Diminishing marginal products
                (negative diagonal) are \emph{not} enough for concavity;
                the cross effect $f_{xy}$ is too strong. Along the diagonal
                the function is $f(t,t)=t^{5/4}$, strictly convex in $t$,
                which refutes concavity directly.
          \item \emph{Level curves and gradient (Vol.~5, Ch.~9, Ch.~10).}
                $x^{3/4}y^{1/2}=c$ gives $y=c^{2}x^{-3/2}$: decreasing convex
                isoquants. $\grad f(1,1)=(\tfrac34,\tfrac12)$, norm
                $\sqrt{\tfrac{9}{16}+\tfrac{4}{16}}=\tfrac{\sqrt{13}}{4}
                \approx0.90$. The isoquant $y=x^{-3/2}$ through $(1,1)$ has
                slope $-\tfrac32$, tangent $(2,-3)$, and
                $\tfrac34\cdot2-\tfrac12\cdot3=0$ \checkmark. The gradient
                points to the upper right, towards higher output, and
                $\sqrt{13}/4$ is the output gain per unit step in that best
                direction.
          \item \emph{Critical points (Vol.~7, Ch.~5).} None in the
                interior, since $f_{x}>0$. The minimum value $0$ is attained
                on the two axes, boundary points of the domain, which the
                first-order condition does not see. (At the corner $(0,0)$
                the gradient from (b) does vanish, but the other minimizers
                on the axes are points where the partials do not even
                exist.)
          \item \emph{Contour sets.} $\{f\ge1\}=\{y\ge x^{-3/2}\}$, the
                region above a convex curve: \textbf{convex}. $\{f\le1\}$ is
                not: $(4,\tfrac18)$ and $(\tfrac14,8)$ lie on $f=1$, but the
                midpoint $(\tfrac{17}{8},\tfrac{65}{16})$ has
                $f>1$ (both coordinates exceed $1$).
          \item \emph{Checklist (Vol.~7, Ch.~3).}
                \emph{Concave: no}, $f(t,t)=t^{5/4}$ is strictly convex.
                \emph{Convex: no}, $f(t,1)=t^{3/4}$ is strictly concave.
                \emph{Quasiconcave: yes}, all upper sets lie above convex
                isoquants (Cobb--Douglas always is).
                \emph{Quasiconvex: no}, by (f).
                \emph{Maximum: no}, unbounded.
                \emph{Minimum: yes}, $0$ on the axes.
                \emph{Increasing: yes}: if $x\le x'$ and $y\le y'$ then
                $x^{3/4}\le x'^{3/4}$ and $y^{1/2}\le y'^{1/2}$, and the
                product of nonnegative factors grows.
                \emph{Decreasing: no}, $f(0,0)=0<1=f(1,1)$.
                \emph{Epigraph convex: no}, $f$ is not convex.
        \end{enumerate}
  \item \textbf{Cobb--Douglas with constant returns.}
        \begin{enumerate}[label=(\alph*),leftmargin=1.8em,itemsep=0.4ex]
          \item \emph{Name (Vol.~5, Ch.~14).} Cobb--Douglas, homogeneous of
                degree $\tfrac13+\tfrac23=1$: \textbf{constant returns to
                scale}. Domain $\R^{2}_{+}$ (closed, convex, unbounded),
                range $[0,\infty)$, since $f(t,t)=t$.
          \item \emph{Derivative and Hessian (Vol.~5, Ch.~10, Ch.~11).}
                On $x,y>0$:
                $f_{x}=\tfrac13x^{-2/3}y^{2/3}$,
                $f_{y}=\tfrac23x^{1/3}y^{-1/3}$, and
                \[
                  \Hess f=\frac29\begin{pmatrix}
                    -x^{-5/3}y^{2/3}& x^{-2/3}y^{-1/3}\\
                    x^{-2/3}y^{-1/3}& -x^{1/3}y^{-4/3}\end{pmatrix},\qquad
                  \det\Hess f=\tfrac{4}{81}x^{-4/3}y^{-2/3}
                  -\tfrac{4}{81}x^{-4/3}y^{-2/3}=0 .
                \]
                Diagonal negative, determinant zero: \textbf{negative
                semidefinite}, not definite. The null direction is the ray
                itself: the first row times $(x,y)$ is
                $\tfrac29\bigl(-x^{-2/3}y^{2/3}+x^{-2/3}y^{2/3}\bigr)=0$,
                and likewise the second. Along a ray $f(tx,ty)=tf(x,y)$ is
                linear, so it has no curvature there.
          \item \emph{Level curves, gradient, Euler (Vol.~5, Ch.~9,
                Ch.~10, Ch.~14).} $x^{1/3}y^{2/3}=c\iff y=c^{3/2}x^{-1/2}$.
                $f(1,8)=1\cdot4=4$, so the isoquant is $y=8x^{-1/2}$, slope
                $-4x^{-3/2}=-4$ at $x=1$. $\grad f(1,8)=
                (\tfrac13\cdot4,\ \tfrac23\cdot\tfrac12)=(\tfrac43,\tfrac13)$,
                norm $\tfrac{\sqrt{17}}{3}\approx1.37$; with tangent
                $(1,-4)$: $\tfrac43-\tfrac43=0$ \checkmark. Steepest ascent
                points mostly in the $x$ direction, because $x$ is the
                scarce factor at $(1,8)$. Euler: $1\cdot\tfrac43+8\cdot
                \tfrac13=4=f(1,8)$ \checkmark.
          \item \emph{Critical points.} None in the interior ($f_{x}>0$);
                the minimum $0$ sits on the axes.
          \item \emph{Contour sets.} $\{f\ge4\}=\{y\ge8x^{-1/2}\}$:
                \textbf{convex}. $\{f\le4\}$: $(16,2)$ and $(\tfrac14,16)$
                lie on $f=4$, the midpoint $(\tfrac{65}{8},9)$ gives
                $f=(\tfrac{65}{8})^{1/3}\cdot9^{2/3}\approx8.7>4$.
                \textbf{Not convex.}
          \item \emph{Checklist (Vol.~7, Ch.~3).}
                \emph{Concave: yes}, the Hessian is negative semidefinite
                on the interior, and concavity carries over to the axes by
                continuity (Cobb--Douglas with $a+b\le1$).
                \emph{Strictly concave: no}, $f(t,t)=t$ is linear along the
                diagonal.
                \emph{Convex: no}, $f(t,1)=t^{1/3}$ is strictly concave.
                \emph{Quasiconcave: yes}, because concave.
                \emph{Quasiconvex: no}, by (e).
                \emph{Maximum: no}. \emph{Minimum: yes}, $0$ on the axes.
                \emph{Increasing: yes}, as in Problem~3.
                \emph{Decreasing: no}.
                \emph{Epigraph convex: no}; the hypograph is convex.
        \end{enumerate}
  \item \textbf{Logarithm of a positive quadratic.}
        \begin{enumerate}[label=(\alph*),leftmargin=1.8em,itemsep=0.4ex]
          \item \emph{Domain and range (Vol.~2, Ch.~1).}
                $q=4+x^{2}+y^{2}\ge4>0$ always, so $D=\R^{2}$. $q$ takes
                every value in $[4,\infty)$ and $\ln$ is increasing, so
                $f(\R^{2})=[\ln4,\infty)$.
          \item \emph{Gradient (Vol.~5, Ch.~10).}
                $\grad f=\Bigl(\dfrac{2x}{q},\dfrac{2y}{q}\Bigr)$,
                continuous everywhere since $q\ne0$: differentiable on all
                of $\R^{2}$.
          \item \emph{Level curves and gradient (Vol.~5, Ch.~9).}
                $f=c\iff x^{2}+y^{2}=\mathrm{e}^{c}-4$: circles about the
                origin for $c>\ln4$, the origin itself for $c=\ln4$, empty
                below. Through $(1,1)$ runs the circle $x^{2}+y^{2}=2$
                (level $\ln6$). $\grad f(1,1)=(\tfrac26,\tfrac26)
                =(\tfrac13,\tfrac13)$, norm $\tfrac{\sqrt2}{3}$; it is a
                multiple of the radius vector $(1,1)$, and a radius is
                orthogonal to a circle: tangent $(1,-1)$,
                $\tfrac13-\tfrac13=0$ \checkmark. It points outwards,
                towards larger circles and higher values.
          \item \emph{Hessian (Vol.~5, Ch.~11).} Quotient rule:
                \[
                  f_{xx}=\frac{2(4-x^{2}+y^{2})}{q^{2}},\quad
                  f_{yy}=\frac{2(4+x^{2}-y^{2})}{q^{2}},\quad
                  f_{xy}=-\frac{4xy}{q^{2}},
                \]
                \[
                  \det\Hess f=\frac{4\bigl[16-(x^{2}-y^{2})^{2}\bigr]
                  -16x^{2}y^{2}}{q^{4}}
                  =\frac{4\bigl[16-(x^{2}+y^{2})^{2}\bigr]}{q^{4}} .
                \]
                At $(0,0)$: $\Hess f=\mathrm{diag}(\tfrac12,\tfrac12)$,
                positive definite. At $(1,1)$ ($q=6$):
                $\begin{pmatrix}2/9&-1/9\\-1/9&2/9\end{pmatrix}$,
                $\det=\tfrac{1}{27}>0$, $f_{xx}>0$: positive definite. At
                $(2,2)$ ($q=12$):
                $\begin{pmatrix}1/18&-1/9\\-1/9&1/18\end{pmatrix}$,
                $\det=-\tfrac{1}{108}<0$: indefinite.
          \item \emph{Critical points (Vol.~7, Ch.~5).} $2x/q=0$ and
                $2y/q=0$ give only $(0,0)$; positive definite Hessian: strict
                local minimum $\ln4$, and global, since $f\ge\ln4$ by (a).
          \item \emph{Contour sets.} $\{f\le\ln8\}=\{x^{2}+y^{2}\le4\}$, a
                disk: \textbf{convex}. $\{f\ge\ln8\}=\{x^{2}+y^{2}\ge4\}$,
                the outside of the disk: it contains $(2,0)$ and $(-2,0)$
                but not their midpoint $(0,0)$. \textbf{Not convex.}
          \item \emph{Checklist (Vol.~7, Ch.~3).}
                \emph{Concave: no}, the Hessian is positive definite at the
                origin.
                \emph{Convex: no}, it is indefinite at $(2,2)$; along the
                $x$-axis $g(t)=\ln(4+t^{2})$ has
                $g''(t)=\dfrac{2(4-t^{2})}{(4+t^{2})^{2}}<0$ for $|t|>2$.
                \emph{Quasiconvex: yes}, every lower set is a disk (or
                empty, or a point).
                \emph{Quasiconcave: no}, by (f).
                \emph{Maximum: no}, unbounded above.
                \emph{Minimum: yes}, $\ln4$ at $(0,0)$.
                \emph{Increasing: no}, $(-1,0)\le(0,0)$ but
                $f(-1,0)=\ln5>\ln4=f(0,0)$.
                \emph{Decreasing: no}, $(0,0)\le(1,0)$ but
                $f(1,0)=\ln5>f(0,0)$.
                \emph{Epigraph convex: no}, $f$ is not convex.
                \emph{Convex on a disk:} the determinant is positive
                exactly when $x^{2}+y^{2}<2^{2}$, and there
                $f_{xx}=2(4-x^{2}+y^{2})/q^{2}>0$ because $x^{2}<4$. So $f$
                is convex on the disk of radius $2$ (closed, by continuity),
                and on no larger disk, since the determinant turns negative
                just outside it. This is the typical shape of a
                quasiconvex but not convex function: a bowl at the bottom
                that flattens into a trumpet.
        \end{enumerate}
\end{solutions}
